\documentclass[14pt,a4paper]{extarticle}

\usepackage[a4paper,left=3cm,right=1cm,top=2cm,bottom=2cm]{geometry}
\usepackage{fontspec}
\usepackage{polyglossia}
\setmainlanguage{russian}
\setotherlanguage{english}
\setmainfont{Times New Roman}

\usepackage{setspace}
\usepackage{array}
\usepackage{longtable}
\usepackage{booktabs}
\usepackage{caption}
\usepackage{enumitem}

\setlength{\parindent}{1.25cm}
\setlength{\parskip}{0pt}
\singlespacing
\pagenumbering{gobble}
\sloppy
\hyphenpenalty=10000
\exhyphenpenalty=10000
\tolerance=1000
\emergencystretch=2em

\captionsetup{justification=centering}
\setlist{nosep}
\renewcommand{\arraystretch}{1.12}

\begin{document}
\begin{titlepage}
  \begin{center}
    Министерство образования Республики Беларусь\\[1em]
    Учреждение образования\\
    БЕЛОРУССКИЙ ГОСУДАРСТВЕННЫЙ УНИВЕРСИТЕТ \\
    ИНФОРМАТИКИ И РАДИОЭЛЕКТРОНИКИ\\[1em]

    \begin{minipage}{\textwidth}
      \begin{flushleft}
        \begin{tabular}{ l l }
          Факультет & информационных технологий и управления\\
          Кафедра   & интеллектуальных информационных технологий
        \end{tabular}
      \end{flushleft}
    \end{minipage}\\[1em]

    {ЗАДАНИЕ}\\
    {по дисциплине «Современная политэкономия»}\\
    {на тему:}\\[1em]
    \textbf{\large ЦИРКУЛЯРНАЯ ЭКОНОМИКА, ЕЁ ОСОБЕННОСТИ В РЕСПУБЛИКЕ БЕЛАРУСЬ}\\[1em]

    \vspace{2em}
    \begin{tabular}{ p{0.65\textwidth}p{0.25\textwidth} }
      Студент гр. 421702 & Н.\,В.~Сечейко \\
      Проверяющий & Е.\,Н.~Макеева \\
    \end{tabular}
    
    \vfill
    {\normalsize Минск 2026}
  \end{center}
\end{titlepage}

\textbf{Циркулярная экономика в Республике Беларусь}

Эта версия сделана как содержательная основа для учебной работы по СПЭ: с определениями, схемами, таблицами, разбором белорусской специфики, кратким резюме стратегии до 2035 года, цифровыми инициативами и подборкой источников за последние 5 лет. Для Беларуси тема уже не теоретическая: в 2024 году правительство утвердило Национальную стратегию развития экономики замкнутого цикла до 2035 года, а в статистике 2024--2026 годов уже видны и опорные достижения, и очень жесткие ограничения системы.

\textbf{Что такое циркулярная экономика}

В рабочем, академически корректном смысле циркулярная экономика --- это не просто переработка мусора. Это модель хозяйствования, в которой ценность материалов, изделий и энергии удерживается в обороте как можно дольше, а образование отходов и потребление первичных ресурсов снижаются. В белорусской стратегии это зафиксировано прямо: циркулярная экономика имеет восстановительный и замкнутый характер, ориентирована на снижение потребления сырьевых и топливно-энергетических ресурсов и на максимальное использование отходов. Международные трактовки у ОЭСР, Европейская комиссия и Ellen MacArthur Foundation совпадают по ядру: нужно устранять отходы на стадии дизайна, продлевать жизнь продуктам и возвращать материалы в цикл использования.

\textbf{Схема. Линейная и циркулярная модели}

Линейная экономика

добыча сырья $\rightarrow$ производство $\rightarrow$ потребление $\rightarrow$ выброс / захоронение

Циркулярная экономика

экодизайн $\rightarrow$ производство $\rightarrow$ потребление $\rightarrow$ повторное использование $\rightarrow$ ремонт / модернизация $\rightarrow$ переработка / рециклинг $\rightarrow$ возврат в производство

Такая схема соответствует белорусской стратегии, где линейный принцип «добывай, производи, выбрасывай» должен сменяться принципом «добывай, производи, повторно используй», а также международной модели циркулярности, где ключевыми считаются продление жизненного цикла, повторное использование и регенерация систем.

\textbf{Схема. Иерархия циркулярных действий}

rethink $\rightarrow$ redesign $\rightarrow$ refuse $\rightarrow$ reduce $\rightarrow$ reuse $\rightarrow$ repurpose $\rightarrow$ repair $\rightarrow$ recycle $\rightarrow$ recover

Это фактически девятизвенная логика белорусской стратегии: от переосмысления бизнес-моделей и улучшения дизайна до переработки и энергетического использования отходов. Важный смысл этой иерархии в том, что переработка --- не первый шаг, а уже поздняя стадия. Сначала нужно не допустить образования отхода, затем продлить жизнь вещи, и только потом перерабатывать.

\textbf{Схема. Какие бизнес-модели относятся к циркулярным}

циркулярные поставки $\rightarrow$ восстановление ресурсов $\rightarrow$ платформы обмена и совместного использования $\rightarrow$ продление жизненного цикла продукции $\rightarrow$ продукт как услуга

Для учебной работы это важно сформулировать отдельно: циркулярная экономика --- это не один сектор отходов, а целый набор бизнес-моделей, который затрагивает дизайн, логистику, аренду, сервис, ремонт, обмен, повторное сырье и цифровые платформы. Белорусская стратегия прямо добавляет: практическая реализация таких моделей требует цифровой поддержки, включая мобильные, облачные, IoT- и робототехнические решения.

\textbf{Отходы как ресурс и инфраструктура переработки в Беларуси}

Если говорить жестко, у Беларуси уже есть база для циркулярной экономики, но пока нет полного замыкания циклов. По официальным данным, в 2024 году в стране образовалось 4 263,0 тыс. т твердых коммунальных отходов, из них использовано 1 686,7 тыс. т, то есть 39,6\%. Это исторический максимум, но одновременно это означает, что примерно 60\% коммунальных отходов все еще не возвращаются в хозяйственный оборот. Отдельно в 2024 году было заготовлено 840,23 тыс. т основных видов вторичных материальных ресурсов: 419,92 тыс. т бумаги и картона, 196,08 тыс. т стекла, 116,56 тыс. т пластмасс, 60,24 тыс. т изношенных шин, 20,91 тыс. т отходов электрического и электронного оборудования и 26,53 тыс. т отработанных масел. Раздельным сбором охвачено 85,3\% населения.

\textbf{Таблица с ключевыми показателями по ТКО и ВМР}

\begin{longtable}{>{\raggedright\arraybackslash}p{4.2cm}>{\raggedright\arraybackslash}p{2.4cm}>{\raggedright\arraybackslash}p{4.0cm}>{\raggedright\arraybackslash}p{2.4cm}}
\toprule
\textbf{Показатель} & \textbf{Значение} & \textbf{Комментарий} & \textbf{Источник} \\
\midrule
\endfirsthead
\toprule
\textbf{Показатель} & \textbf{Значение} & \textbf{Комментарий} & \textbf{Источник} \\
\midrule
\endhead
Образование ТКО в 2024 году & 4 263,0 тыс. т & Национальный масштаб потока коммунальных отходов & Экологический бюллетень \\
Использование ТКО в 2024 году & 1 686,7 тыс. т & Максимум за весь период наблюдений & Экологический бюллетень \\
Уровень использования ТКО & 39,6\% & Существенный рост, но далеко до цели 90\% & Экологический бюллетень; госпрограмма 2026--2030 \\
Захоронение ТКО & 2 576,3 тыс. т & Значительная доля все еще идет на полигоны & Экологический бюллетень \\
Заготовлено основных ВМР & 840,23 тыс. т & Сырьевая база для переработчиков & Экологический бюллетень \\
Бумага и картон & 419,92 тыс. т & Крупнейшая фракция ВМР & Экологический бюллетень \\
Стекло & 196,08 тыс. т & Вторая по объему фракция & Экологический бюллетень \\
Пластмассы & 116,56 тыс. т & Критически важная фракция для роста переработки & Экологический бюллетень \\
Изношенные шины & 60,24 тыс. т & Поток для крошки, RDF и других применений & Экологический бюллетень \\
Охват населения раздельным сбором & 85,3\% & Сильная сторона системы & Госпрограмма 2026--2030 \\
\bottomrule
\end{longtable}

Здесь важно не перепутать даты и не попасть в ловушку противоречивых цифр. В документах за 2025 год фигурирует 9 мусороперерабатывающих заводов и 81 линия сортировки с мощностью около 1,1 млн т ТКО в год, а в госпрограмме на 2026--2030 годы --- уже 10 заводов и 81 линия, которые позволяют сортировать около 1,8 млн т в год. Это не ошибка, а следствие ввода новых мощностей в 2025 году, прежде всего в Минск, Бобруйск и Новополоцк. В учебном тексте это лучше прямо оговорить: цифры нужно сопровождать пометкой «по состоянию на 2025 год» или «с учетом вводов 2025 года».

Еще одна принципиальная вещь: большая часть белорусских объектов на первой стадии не «глубоко перерабатывает», а сортирует и готовит вторсырье --- разделяет по видам, маркам, прессует, дробит и отправляет профильным переработчикам. Это системно важно, потому что сортировка сама по себе не равна циркулярности: если нет следующего звена переработки и сбыта вторичного сырья, то значимая часть потоков уходит в захоронение. Именно поэтому в официальных материалах отдельно отмечаются дефицит качественной макулатуры и ограниченные мощности по переработке нетарного стеклобоя.

\textbf{Схема. Как отход превращается в ресурс в белорусской системе}

домохозяйства / бизнес $\rightarrow$ раздельный сбор и смешанный поток ТКО $\rightarrow$ сортировка $\rightarrow$ извлечение ВМР $\rightarrow$ первичная подготовка (прессование, дробление, брикетирование) $\rightarrow$ профильный переработчик $\rightarrow$ новая продукция / RDF / компост / биотопливо

Для Беларуси эта схема уже работает, но не полностью. Бумага, стекло, полимеры, шины и электронные отходы возвращаются в экономику лучше, чем органика, смешанные упаковки и часть строительных или крупнотоннажных отходов. Именно здесь сегодня главный разрыв между «сбором» и настоящим замыканием цикла.

\textbf{Таблица с ключевыми объектами сортировки и переработки ТКО}

\begin{longtable}{>{\raggedright\arraybackslash}p{5.55cm}>{\raggedright\arraybackslash}p{10.0cm}}
\toprule
\textbf{Объект} & \textbf{Данные} \\
\midrule
\endfirsthead
\toprule
\textbf{Объект} & \textbf{Данные} \\
\midrule
\endhead
Брестский мусороперерабатывающий завод & \textbf{Расположение:} Брест. \textbf{Специализация:} региональная сортировка ТКО, извлечение ВМР, работа по брестской зоне обслуживания. \textbf{Объем / мощность:} фактический поток около 103 тыс. т отходов в год; обслуживает Брест и соседние районы. \textbf{Источник:} БЕЛТА 2021; обзор МЖКХ 2025. \\
Комплекс сортировки ТКО в Барановичах & \textbf{Расположение:} Барановичи. \textbf{Специализация:} сортировка ТКО, компостирование, подготовка RDF. \textbf{Объем / мощность:} 75 тыс. т/год; в проектных материалах --- 74,8 тыс. т/год. \textbf{Источник:} БЕЛТА; материалы общественного обсуждения. \\
Оршанский региональный комплекс по обращению с ТКО & \textbf{Расположение:} Орша. \textbf{Специализация:} сортировка ТКО, извлечение ВМР, получение пре-RDF, вторичный щебень, щепа. \textbf{Объем / мощность:} 60 тыс. т/год. \textbf{Источник:} БЕЛТА; оператор ВМР. \\
Гродненский завод по утилизации и механической сортировке отходов & \textbf{Расположение:} Гродно. \textbf{Специализация:} полный цикл сортировки ТКО, производство RDF. \textbf{Объем / мощность:} до 103 тыс. т отходов в год; один из первых заводов с полным циклом и RDF. \textbf{Источник:} Grodno News; БЕЛТА. \\
Витебский мусоросортировочный завод & \textbf{Расположение:} Витебск. \textbf{Специализация:} сортировка ТКО, извлечение бумаги, стекла, полимеров, шин. \textbf{Объем / мощность:} проектная мощность 100 тыс. т/год; фактически через завод в 2024 году прошло 47 216 т ТКО. \textbf{Источник:} Витьбичи; Витебский горисполком. \\
Могилевский мусороперерабатывающий завод & \textbf{Расположение:} Могилев. \textbf{Специализация:} сортировка ТКО и выделение вторичных ресурсов. \textbf{Объем / мощность:} 70 тыс. т/год. \textbf{Источник:} Могилевский горисполком. \\
Минский мусороперерабатывающий завод & \textbf{Расположение:} Минск. \textbf{Специализация:} новый крупный комплекс по обращению с ТКО, автоматическая сортировка, RDF, перспективное компостирование. \textbf{Объем / мощность:} 600 тыс. т/год; план выпуска около 113 тыс. т RDF; при вводе второй очереди заявлен уровень переработки до 90\%. \textbf{Источник:} БЕЛТА; CTV; Minsk News. \\
Новополоцкий региональный комплекс по обращению с ТКО & \textbf{Расположение:} Новополоцк. \textbf{Специализация:} сортировка ТКО, переработка строительных и древесных отходов, утилизация био- и медотходов, гранулы из полиэтилена. \textbf{Объем / мощность:} 80 тыс. т ТКО/год, 10 тыс. т строительных отходов, 4 тыс. т зеленой массы, 700 т био- и медотходов, 500 т вторичных гранул. \textbf{Источник:} Витебский облкомприроды; БЕЛТА. \\
Бобруйский мусороперерабатывающий завод & \textbf{Расположение:} Бобруйский район. \textbf{Специализация:} сортировка ТКО, переработка полиэтилена и пластика в флексу и гранулы. \textbf{Объем / мощность:} поток зоны к 2025 году вырос до 60 тыс. т/год плюс отходы района; завод рассчитан на глубокую работу с пластиком и ПЭТ. \textbf{Источник:} Коммерческий курьер. \\
Гомельский региональный комплекс по обращению с ТКО & \textbf{Расположение:} Гомель. \textbf{Специализация:} крупный региональный комплекс, сортировка ТКО и рост изъятия ВМР. \textbf{Объем / мощность:} на стадии строительства; официально заявлена мощность 230 тыс. т/год. \textbf{Источник:} Гомельский облкомприроды; БЕЛТА. \\
\bottomrule
\end{longtable}

Из этой таблицы видно главное: система в Беларуси движется от локальных районных линий к региональным межрайонным комплексам. Это правильно с точки зрения экономики масштаба. Но второй, более неприятный вывод такой: даже при заметном росте сортировочных мощностей страна все еще упирается в недостаток последующей глубокой переработки по части фракций, а значит, сама по себе стройка новых комплексов не гарантирует замыкания цикла.

\textbf{Плюсы и ограничения для Беларуси}

Для Беларуси преимущества циркулярной экономики не абстрактные, а вполне прикладные. Во-первых, это снижение ресурсо- и энергоемкости: стратегия прямо ставит цель уменьшить энергоемкость ВВП, снизить интенсивность образования отходов производства и повысить использование ТКО и отходов производства. То есть циркулярность рассматривается не как «экологическая опция», а как способ повысить продуктивность ресурсов. В международной логике это совпадает с подходом Программа ООН по окружающей среде и ОЭСР: меньше зависимость от первичного сырья и меньше экологическая нагрузка на единицу выпуска.

Во-вторых, Беларусь уже имеет несколько сильных стартовых позиций. Традиционные виды вторичных ресурсов по отходам производства --- макулатура, стеклобой, полимерные отходы, изношенные шины --- в 2024 году использовались более чем на 90\%. Раздельный сбор охватывает большую часть населения. Страна также развивает новые направления: компостирование органической части ТКО и производство RDF-топлива для цементной промышленности. Это означает, что базовый фундамент для масштабирования не нужно строить с нуля.

Во-третьих, у циркулярной модели есть и промышленная выгода. Там, где отходы становятся сырьем или топливом, сокращаются расходы на импортируемые материалы и энергоносители, формируются новые цепочки добавленной стоимости, а также растет спрос на оборудование, сервис, логистику, цифровой мониторинг и инженерные решения. Для страны с сильной ролью промышленности, ЖКХ, строительства и аграрного сектора это особенно важно, потому что стратегия называет именно эти отрасли приоритетными для циркулярной трансформации.

Но слабые места у Беларуси не менее очевидны. Главный системный минус --- разрыв между сбором, сортировкой и реальным конечным использованием. В официальных материалах прямо сказано, что организации, осуществляющие сортировку, как правило, не перерабатывают ВМР глубоко, а только готовят их к дальнейшему использованию. Более того, часть потенциально полезных отходов все еще идет на захоронение из-за отсутствия или недостаточности перерабатывающих мощностей. Это и есть самое болезненное место всей конструкции. Если не закрыть его, система будет производить не «циркулярность», а просто более дорогую сортировку.

Второй большой барьер --- структура промышленных отходов. В 2024 году около 97,5\% накопленных отходов производства приходилось на крупнотоннажные потоки --- галитовые отходы, глинисто-солевые шламы и фосфогипс. По ним циркулярность пока очень ограничена: ежегодно на хранение уходит около 35 млн т таких отходов. Это означает, что красивая статистика по бумаге, стеклу и шинам не должна вводить в заблуждение: самый тяжелый по массе сегмент отходов все еще остается слабым местом.

Третий барьер --- инфраструктурное наследие. В стратегических экологических документах подчеркивается, что более 90\% действующих полигонов строились еще в советский период, а проектные мощности многих объектов практически исчерпаны. То есть страна одновременно решает две задачи: ей надо и наращивать сортировку/переработку, и просто успевать обновлять базовую инфраструктуру захоронения и экологической защиты. Это повышает капиталоемкость перехода и делает процесс медленнее, чем в теории.

Жесткий вывод здесь такой: для Беларуси циркулярная экономика выгодна и нужна, но самый опасный самообман --- считать, что контейнеры, линии сортировки и новые заводы сами по себе уже означают замкнутый цикл. Нет. Настоящий результат начинается там, где у каждой собранной фракции есть стабильный downstream: переработчик, рынок сбыта, технический стандарт, логистика и цифровая прослеживаемость. Без этого цель 90\% по ТКО останется нормативной, а не фактической.

\textbf{Национальная стратегия до 2035 года}

Базовый документ для темы --- постановление Совета Министров № 393 от 29 мая 2024 года, официально опубликованное на Национальный правовой Интернет-портал Республики Беларусь 6 июня 2024 года. Документ задает не «мусорную» политику в узком смысле, а рамку перехода всей экономики к более замкнутому использованию ресурсов. В нем прямо названы приоритетные направления: экодизайн, ресурсоэффективное производство, промышленный симбиоз, сфера упаковки и шеринговая экономика; приоритетные сектора --- промышленность, строительство, ЖКХ, транспорт, сельское и лесное хозяйство, сфера услуг.

\textbf{Схема. Логика Национальной стратегии}

приоритеты

экодизайн + ресурсоэффективность + промышленный симбиоз + упаковка + sharing

сектора

промышленность + строительство + ЖКХ + транспорт + сельское и лесное хозяйство + услуги

инструменты

планирование + НДТМ + мониторинг + инвестиционные проекты + обучение кадров + форумы и популяризация + информационная платформа

результат

меньше первичного сырья, меньше отходов, выше использование ТКО и отходов производства, ниже энергоемкость, выше ресурсоэффективность

Эта внутренняя архитектура стратегии хорошо показывает ее сильную сторону: документ не замыкается на одном министерстве и не сводит тему только к переработке. Но в этом же и ее риск: чем шире рамка, тем больше шанс, что часть задач будет «размазана» по ведомствам без реального сквозного центра принятия решений.

\textbf{Краткое резюме стратегии}

Стратегия задает для Беларуси переход от линейной к восстановительной модели использования ресурсов. Она признает, что циркулярность --- это одновременно производственная, технологическая, институциональная и поведенческая трансформация. Документ привязывает циркулярную экономику не только к отходам, но и к конкурентоспособности, ресурсной производительности, инновациям, климатической и экологической повестке. Отдельно важно, что в приложение вынесен план мероприятий: создание межведомственной рабочей группы, включение циркулярных мер в государственное и отраслевое планирование, внедрение наилучших доступных технических методов, формирование системы мониторинга, инвестиционного портфеля, образовательных материалов и платформы общего доступа к результатам мониторинга.

Самое полезное для учебной работы --- не переписывать стратегию по главам, а показать ее в пяти тезисах.

Во-первых, Беларусь официально признала циркулярную экономику отдельным направлением долгосрочной экономической политики.

Во-вторых, приоритет поставлен не на ликвидацию последствий, а на предотвращение образования отходов, экодизайн и ресурсоэффективность.

В-третьих, переход должен идти через приоритетные сектора с учетом отраслевой специфики.

В-четвертых, стратегия требует цифрового мониторинга и информационной платформы, то есть сразу предполагает data-driven управление.

В-пятых, к 2035 году заданы конкретные измеримые индикаторы, по которым стратегию можно проверять не лозунгами, а цифрами.

\textbf{Таблица с целевыми индикаторами стратегии на 2035 год}

\begin{longtable}{>{\raggedright\arraybackslash}p{6.0cm}>{\raggedright\arraybackslash}p{3.2cm}>{\raggedright\arraybackslash}p{3.2cm}}
\toprule
\textbf{Индикатор} & \textbf{База} & \textbf{Цель 2035} \\
\midrule
\endfirsthead
\toprule
\textbf{Индикатор} & \textbf{База} & \textbf{Цель 2035} \\
\midrule
\endhead
Энергоемкость ВВП, кг у.т./млн руб. & 388,1 в 2021 г. & 268 \\
Доля первичной энергии из ВИЭ в валовом потреблении ТЭР, \% & 7,9 в 2021 г. & 9 \\
Интенсивность образования отходов производства на единицу ВВП, кг/руб. & 0,49 в 2021 г. & 0,40 \\
Использование отходов производства без крупнотоннажных, \% & --- & не менее 90 \\
Использование ТКО в общем объеме их образования, \% & 31,1 в 2021 г. & 90 \\
Эффективность водопользования, руб./м$^3$ & 60,7 в 2021 г. & 62 \\
\bottomrule
\end{longtable}

Источник: Национальная стратегия развития экономики замкнутого цикла до 2035 года.

Самая амбициозная и одновременно самая уязвимая цель здесь --- рост использования ТКО с 31,1\% до 90\%. Это огромный разрыв. Он означает, что Беларусь должна не просто улучшить сортировку, а радикально нарастить переработку, работу с органикой, сбыт вторичного сырья, производство RDF и управленческую прозрачность потоков отходов. Если честно, это уже не «экологическая донастройка», а масштабная перестройка инфраструктуры и рынков.

\textbf{Цифровые инициативы}

Цифровой слой в циркулярной экономике --- не украшение, а управленческий минимум. Без данных о составе, качестве, происхождении, маршруте и остаточной ценности материалов никакой замкнутый цикл не работает. Это прямо признают и белорусская стратегия, и аналитика ОЭСР: цифровизация нужна для дизайна, учета потоков, мониторинга, обмена ресурсами и enforcement.

\textbf{Зарубежный опыт}

\begin{longtable}{>{\raggedright\arraybackslash}p{2.8cm}>{\raggedright\arraybackslash}p{3.2cm}>{\raggedright\arraybackslash}p{4.2cm}>{\raggedright\arraybackslash}p{2.2cm}}
\toprule
\textbf{Инициатива} & \textbf{Что это такое} & \textbf{Для чего нужно в циркулярной экономике} & \textbf{Источник} \\
\midrule
\endfirsthead
\toprule
\textbf{Инициатива} & \textbf{Что это такое} & \textbf{Для чего нужно в циркулярной экономике} & \textbf{Источник} \\
\midrule
\endhead
Цифровой паспорт продукта и проект CIRPASS & Система цифровой информации о составе, происхождении, ремонте и переработке продукта & Делает продукт прослеживаемым в течение жизненного цикла и поддерживает повторное использование, ремонт и рециклинг & Европейская комиссия / ECESP / CIRPASS \\
Madaster & Онлайн-реестр материалов для зданий и инфраструктуры & Позволяет заранее знать, какие материалы можно разобрать и повторно использовать в строительстве & ECESP; Madaster \\
Excess Materials Exchange & B2B-платформа обмена избыточными материалами & Превращает «лишние» материалы из отхода в актив и автоматически ищет им следующее применение & Holland Circular Hotspot; платформа EME \\
AMP Robotics & ИИ-платформа компьютерного зрения для сортировки отходов & Повышает чистоту фракций и извлечение ценных материалов на сортировке & Ellen MacArthur Foundation; AMP \\
\bottomrule
\end{longtable}

Из зарубежного опыта для Беларуси особенно важны не «красивые кейсы», а четыре прикладные идеи. Первая --- паспортизация продукта или материала. Вторая --- цифровые биржи и маркетплейсы вторичных ресурсов. Третья --- ИИ-сортировка с компьютерным зрением. Четвертая --- цифровая отчетность и контроль соответствия. Именно они создают мост между сбором отходов и их реальным возвращением в экономику.

\textbf{Белорусская практика и локальные разработки}

\begin{longtable}{>{\raggedright\arraybackslash}p{5.0cm}>{\raggedright\arraybackslash}p{10.6cm}}
\toprule
\textbf{Инициатива} & \textbf{Данные} \\
\midrule
\endfirsthead
\toprule
\textbf{Инициатива} & \textbf{Данные} \\
\midrule
\endhead
Оператор вторичных материальных ресурсов: карта приемных пунктов и реестр & \textbf{Формат:} национальный сайт с картой и реестром организаций. \textbf{Практическая роль:} дает пользователю и бизнесу доступ к сети из более чем 1630 пунктов приема вторсырья и к официальной инфраструктуре системы. \textbf{Источник:} Оператор ВМР. \\
Портал и приложение 115.бел & \textbf{Формат:} государственный цифровой сервис ЖКХ. \textbf{Практическая роль:} позволяет жителям фиксировать проблемы, в том числе мусор, свалки, загрязнение территории, с фото и геопривязкой. \textbf{Источник:} 115.бел; App Store / Google Play. \\
Белорусская универсальная товарная биржа: электронные торги неликвидами & \textbf{Формат:} B2B-площадка. \textbf{Практическая роль:} помогает превращать невостребованные остатки и неликвидные активы в оборотный ресурс, что по сути работает как элемент промышленного симбиоза. \textbf{Источник:} БУТБ; БЕЛТА. \\
Ecoplatform в Минске & \textbf{Формат:} сеть фандоматов с приложением и бонусами. \textbf{Практическая роль:} стимулирует население сдавать ПЭТ и алюминий через финансовую мотивацию и геймификацию. \textbf{Источник:} Belretail; Minsk News; Google Play. \\
Оптические сепараторы на новых комплексах & \textbf{Формат:} элемент цифровизированной сортировки. \textbf{Практическая роль:} позволяют распознавать материал по форме, цвету и структуре и повышать качество ВМР. \textbf{Источник:} БЕЛТА; Витьбичи. \\
\bottomrule
\end{longtable}

Здесь нужна честная формулировка. В Беларуси уже есть цифровые элементы циркулярной инфраструктуры, но пока нет единой национальной цифровой экосистемы циркулярной экономики. Есть карта пунктов приема, реестр операторов, гражданский сервис фиксации проблем, биржевая площадка для неликвидов, локальные reverse-vending решения и частичная цифровизация сортировочных комплексов. Но нет полноценных material passports, нет национального цифрового обмена промышленными побочными потоками и нет сквозной публичной аналитики по составу и качеству отходов в реальном времени. Именно это и отличает стадию «становления» от зрелой циркулярной системы.

\textbf{Перспективы и ИИ}

Если смотреть на перспективы без розовых очков, то ключевая задача Беларуси до 2035 года --- не просто строить еще линии сортировки, а перейти от инфраструктуры сбора к инфраструктуре принятия решений. Стратегия уже требует систему мониторинга и информационную платформу общего доступа. Это правильный сигнал: следующий этап --- не контейнеры как таковые, а данные о потоках, загрязненности, логистике, качестве вторсырья, остаточном ресурсе изделий и экономике повторного использования.

\textbf{Схема. Где ИИ и цифровые сервисы действительно полезны}

сбор данных $\rightarrow$ анализ состава отходов $\rightarrow$ прогноз потоков $\rightarrow$ автоматическая сортировка $\rightarrow$ поиск следующего применения материала $\rightarrow$ контроль выполнения целей стратегии

В этой цепочке ИИ и цифровые инструменты могут работать на каждом шаге: компьютерное зрение распознает фракции на конвейере, аналитика предсказывает сезонные пики отходов, платформы связывают продавца вторсырья с переработчиком, чат-боты помогают населению правильно сортировать, а цифровые паспорта материалов снижают потери при ремонте и демонтаже. Именно такое понимание цифровизации дают ОЭСР, европейские DPP-проекты и кейсы ИИ-сортировки.

Для Беларуси наиболее реалистичный набор перспективных решений выглядит так.

Во-первых, единое мобильное приложение циркулярной навигации для граждан: что и куда сдавать, где ближайший пункт или фандомат, как подготовить отход, когда вывозят раздельно собранные фракции, какие бонусы доступны, какие ошибки чаще всего совершают жители. Сейчас элементы этого уже разрозненно существуют, но не собраны в единую экосистему.

Во-вторых, B2B-платформа промышленного симбиоза. Белорусская стратегия прямо называет промышленный симбиоз одним из приоритетов. Значит, логичный следующий шаг --- цифровая площадка, где предприятие может публиковать состав, объем, периодичность и качество побочного потока, а другое предприятие --- находить его как сырье. Это был бы уже не просто «рынок отходов», а цифровая инфраструктура для замыкания производственных циклов. Это вывод, вытекающий из сочетания самой стратегии и зарубежных обменных платформ вроде Excess Materials Exchange.

В-третьих, цифровые паспорта для строительных материалов и отходов строительства. Для Беларуси это особенно перспективно, потому что строительство --- один из приоритетных секторов стратегии, а строительные отходы уже начинают обрабатываться на новых региональных комплексах, например в Новополоцке. Паспортизация материалов позволила бы заранее планировать демонтаж, повторное использование и переработку. Это не фантазия, а логическое продолжение модели Madaster, адаптированной под национальный рынок.

В-четвертых, ИИ-мониторинг сортировочных комплексов. Новые и модернизированные белорусские комплексы уже используют оптические сепараторы. Следующий шаг --- накопление датасетов по загрязненности фракций, распознаванию материалов и качеству потока, а затем --- обучение моделей, которые будут не просто сортировать, а подсказывать оптимальную скорость линии, настройки отбора и приоритетные маршруты для конкретной смеси отходов. Это уже не теория: мировая практика это применяет, а белорусская инфраструктура частично готова к такой надстройке.

В-пятых, чат-боты и цифровые помощники для РОП, отчетности и сортировки. Белорусская стратегия предусматривает мониторинг, обучение и информационную платформу, а в системе обращения с отходами уже есть официальный реестр и карта приемных пунктов. Значит, вполне реалистично создать прикладных ботов для бизнеса и населения: для разъяснения требований расширенной ответственности производителя, поиска канала сдачи отходов, подсказок по маркировке упаковки, проверки, перерабатывается ли материал, и формирования первичной отчетности. Это как раз тот случай, где ИИ дешевле и полезнее, чем очередная бумажная памятка.

Самый трезвый прогноз такой: у Беларуси есть шанс быстро продвинуться в циркулярной экономике не за счет одной «мегатехнологии», а за счет комбинации региональных комплексов + цифрового мониторинга + рынков вторсырья + работы с органикой и RDF + нормальной отраслевой аналитики. Если этого не произойдет, страна рискует попасть в типичную ловушку переходного периода: много капитальных объектов, много отчетности, но слишком мало реально замкнутых циклов.

\textbf{Источники за 2021--2026 годы}

Ниже --- подборка, которую уже можно использовать как основу для списка литературы: официальные электронные источники для фактов по Беларуси, монографии и учебные издания для теории, а также свежие конференционные статьи для современной научной дискуссии.

\textbf{Электронные источники}

О Национальной стратегии развития экономики замкнутого цикла Республики Беларусь на период до 2035 года: постановление Совета Министров Республики Беларусь от 29.05.2024 № 393. 2024.

Государственная программа «Экология» на 2026--2030 годы. 2026.

Состояние природной среды Беларуси. Экологический бюллетень. Обращение с отходами в 2024 году. 2025.

Материалы Министерство природных ресурсов и охраны окружающей среды Республики Беларусь о развитии системы обращения с отходами, региональных комплексах и стратегических документах. 2024--2026.

Материалы Министерство жилищно-коммунального хозяйства Республики Беларусь и Оператора ВМР о региональных комплексах, сети приема и финансировании инфраструктуры. 2025--2026.

OECD. Digitalisation for the Transition to a Resource Efficient and Circular Economy. 2024.

European Circular Economy Stakeholder Platform. Digital Product Passport as Enabler for the Circular Economy. 2023.

UNEP. Circularity and related circular economy resources. 2024--2025.

\textbf{Монографии}

Ильина, Е. А. Циркулярная экономика: концептуальные подходы и механизмы реализации. Монография. Минск: IBB, 2021. ISBN 978-985-7261-22-2.

Аронов, И. З., Еремеева, Н. В., Панюкова, В. В. и др. Экономика замкнутого цикла и техническое регулирование. Монография. Москва: МГИМО-Университет, 2024. 644 с. ISBN 978-5-9228-2823-9.

\textbf{Учебные издания}

Валько, Д. В. Основы циркулярной экономики: учебное пособие для вузов. Москва: Юрайт, 2024. 104 с. ISBN 978-5-534-18659-8.

Рязанова, О. Е., Золотарева, В. П. Циркулярная экономика: учебное пособие. Москва: КноРус, 2022. 181 с. ISBN 978-5-406-09864-6.

Рязанова, О. Е., Золотарева, В. П. Циркулярная экономика (экономика замкнутого цикла) (с практикумом): учебник. Москва: КноРус, 2025. 278 с. ISBN 978-5-406-12738-4.

\textbf{Статьи по результатам конференций}

Батова, Н. Н. Развитие циркулярной экономики в Республике Беларусь: от пилотных проектов к обществу замкнутого цикла. В: современные материалы БГУ, 2024.

Цедрик, А. В. Научный дискурс по проекту постановления «О Национальной стратегии развития экономики замкнутого цикла (циркулярной экономики) Республики Беларусь на период до 2035 года». 2024.

Костюченко, Е. А., Тихолаз, А. В. Циркулярная экономика и проблемы оценки эффективности рециклинга. В: Устойчивое развитие экономики: международные и национальные аспекты, VI Международная научно-практическая конференция, Новополоцк, 31.10--01.11.2024. Изд. 2025.

Головач, О. В. Экономика замкнутого цикла и развитие ее информационного учетно-аналитического обеспечения. В: Устойчивое развитие экономики: международные и национальные аспекты, VI Международная научно-практическая конференция, Новополоцк, 31.10--01.11.2024. Изд. 2025.

Курьян, Е. В., Якуш, Е. О., Гришаева, К. А. Циркулярная экономика: теория и практика. В: материалы международной научно-практической конференции, Новополоцк, 2025.

\end{document}
